\documentclass[../Main.tex]{subfiles}
\begin{document}

\section{Further Work}
\label{chap:further_work}

The work we have done in this specialization project will serve as the foundation for our master’s thesis. We will use the next period to develop, test, and analyze the concept in order to evaluate the concept and its potential. The next step is to create a detailed implementation plan including elements such as developer consideration, software architecture description, feature prioritization, and more. It is also crucial to settle how we want to include the game design elements mentioned in Chapter \ref{chap:game_design}. The prototype that we created during this past period is a great starting point, but naturally, a more detailed and structured plan is needed before continuing with the development. The prototype was mainly developed for testing the central aspects of technology. During the planning phase, we may encounter new ideas or improvements, and thus may need to tweak the idea to some degree.

After the planning phase is completed, we will transition to the development phase where the focus will be to realize the concept in the best possible way. An important measure during this work is to get continuous feedback from users allowing for multiple iterations. The plan that we initially proposed may not work as well as expected or the game design elements may not have the intended effect on the player. As a result, the development phase will be based on agile principles.

When we feel confident that the application is developed such that the main ideas of Hover are realized in a proper way, the real user testing will begin. The goal of the project is to evaluate the potential of the concept and thus the application should focus on player experience and not necessarily polishing the user interface. The most essential consideration is the ability to include all the fundamental game design elements. In addition to this, analytical tools should be implemented in the game. They are crucial for collecting concrete usage data during the testing period. Questions such as “How often does the player open the game?” or “How long does each play session last?” can easily be answered using analytical tools. The testing period will look at the usage and habits of players’ in the game.

Feedback and results from the testing period will be thoroughly analyzed and discussed. We will try to find statistically significant differences that can help us evaluate the concept in a good way. Based on the results and evaluation, a conclusion of how well the concept was received will be drawn. The conclusion should also answer other defined research questions.

\biblio % Needed for referencing to working when compiling individual subfiles - Do not remove
\end{document}